\documentclass[11pt]{article}
\usepackage[margin=1in]{geometry}
\usepackage{amsmath,amssymb,amsfonts,amsthm,mathtools,bm}
\usepackage[T1]{fontenc}
\usepackage{microtype}
\newtheorem{theorem}{Theorem}
\newtheorem{lemma}{Lemma}
\newtheorem{cor}{Corollary}
\newtheorem{defn}{Definition}
\newtheorem{principle}{Principle}
\newtheorem{statement}{Statement}
\theoremstyle{remark}
\newtheorem{statementr}{Statement}

\title{Geometrically Formulating Bohmian Theories on Hermitian Vector Bundles}
\author{James Taylor}
\date{Working draft derived from the 2003 thesis}
\begin{document}
\maketitle

\section{Geometrically formulating Bohmian
  theories} \label{c:geom} 


In this chapter, we shall formulate a Bohmian dynamics defined on a general
Riemannian manifold  \ensuremath{\mathcal{Q}}\ with the wave function\ being a section of a 
cc-Hermitian
vector bundle $E$ over \ensuremath{\mathcal{Q}}. 
 We begin with scalar wave functions over \ensuremath{\mathcal{Q}}\ and then 
 generalize 
to vector bundles. This sets up the general 
non-relativistic 
Bohmian framework. The Pauli equation may be viewed as both a special
case as well as a generalization of this
setup.
We explain how to view the equation as expressing a generalized
version of the Laplacian. After making a few standard observations, 
we finish this chapter with a discussion of how a desire for an
ultralocal Bohmian theory  leads to the Dirac equation. 
It is rooted in a structure which we call the square-root of
the tangent space. 
It is intimately related to the role of Clifford algebra in the Dirac
equation. 

\section{A quick review of tensor analysis}

In this section, we shall give a review of the uses and abuses of
tensor analysis. An appreciation of these methods is essential for an
understanding of this chapter. For most of this section and the rest
of the chapter, the statements can either be thought of as in terms of
the whole bundle or just at a point. Even though it is best to think
of all the structures as just at a point, we shall write our notations
without denoting the point. For example, we shall write
$T\ensuremath{\mathcal{Q}}$ instead of $T_q \ensuremath{\mathcal{Q}}$ even though we will actually be 
discussing the tangent space at the point $q$. 

We start with a metric $g$ on the manifold \ensuremath{\mathcal{Q}}. It is a 
$(0,2)$-tensor meaning that it is an element of $(\ensuremath{T^{*} \ensuremath{\mathcal{Q}}})\ensuremath{^{\otimes 2} }$. It acts
on $(\ensuremath{T \ensuremath{\mathcal{Q}}})\ensuremath{^{\otimes 2} }$; alternatively one could say that it is a bilinear form
acting on tangent vectors. In coordinates, $ g = g_{\ensuremath{j} \ensuremath{k}}
dx^{\ensuremath{j}} \otimes dx^{\ensuremath{k}}$ where the summation convention on
repeated indices is implied. An important use of $g$ is to lower an index,
i.e.\ $g$ may be viewed as a linear map  from $\ensuremath{T \ensuremath{\mathcal{Q}}}$ to $\ensuremath{T^{*} \ensuremath{\mathcal{Q}}}$. Indeed,
given a tangent vector $v$, $g(v, \cdot)$ is a 1-form. If it is
necessary, we will denote the role of $g$ as a loweing operator by
writing $g_{\flat}$. In other words, $g_{\flat}(v) \ensuremath{:=} g(v,
\cdot)$. Even more compactly, one may write $v_{\flat}$ indicating
that $v$ has been converted to a 1-form using $g$.

An assumption
on a metric is that it is nondegenerate which means that given a
vector $v$, there is a $w$ such that $g(v,w) \neq 0$. This assumption
is equivalent to demanding that $g$ as a lowering operator has an
inverse. We denote the raising operator as $G^{\sharp}$ and it is a map from
\ensuremath{T^{*} \ensuremath{\mathcal{Q}}}\ to \ensuremath{T \ensuremath{\mathcal{Q}}}. One can then define $G$ as a bilinear form acting on
$\ensuremath{T^{*} \ensuremath{\mathcal{Q}}}$ by defining $G(\omega, \nu) \ensuremath{:=} \omega(G^{\sharp}(\nu))$ where
$\omega$ and $\nu$ are 1-forms. Such a bilinear form is the same as a
$(2,0)$-tensor i.e.\ it is an elment of $(\ensuremath{T \ensuremath{\mathcal{Q}}}) \ensuremath{^{\otimes 2} }$. In coordinates,
$G = G^{\ensuremath{j} \ensuremath{k}} \partial_{\ensuremath{j}} \otimes \partial_{\ensuremath{k}}$
where $(G^{\ensuremath{j} \ensuremath{k}})$ is the inverse matrix to $(g_{\ensuremath{j}
  \ensuremath{k}})$. 

Most of our use of $G$ is to contract it with $(0,2)$-tensors. On
product elements, $G(\omega \otimes \nu) \ensuremath{:=} G(\omega, \nu)$; it is
then extended by linearity to all of $\ensuremath{T^{*} \ensuremath{\mathcal{Q}}}\ensuremath{^{\otimes 2} }$. Let $\beta$ be a
$(0,2)$-tensor. Then the notations that may be used are $G(\beta) = G
\circ \beta = \beta(G) = G^{\ensuremath{j} \ensuremath{k}} \beta_{\ensuremath{j} \ensuremath{k}}$.    


We introduce some more notation. The
symmetric subspace of a $i$-fold tensor product of a vector space $W$
with itself is denoted by $\ensuremath{\bigodot^{i}{W}}$. The antisymmetric part of such
a tensor product is $\ensuremath{\bigwedge^{i}{W}}$. The direct sum over $i$ of the antisymmetric
parts is the exterior algebra where the product is the wedge product
which may be defined as the antisymmetrization of the tensor product
of the terms. 


A metric is a symmetric tensor, i.e. $g \in \ensuremath{\bigodot^{2}{\ensuremath{T^{*} \ensuremath{\mathcal{Q}}}}}$. Its dual,
$G$, is also symmetric. This means that $G$, when acting on a $(0,2)$
tensor, will only notice the symmetric part of the tensor. In
particular, if $\beta \in \ensuremath{\bigwedge^{2}{\ensuremath{T^{*} \ensuremath{\mathcal{Q}}}}}$, then $G(\beta) = 0$. 

The final element of a metric on a manifold is that it should be
parallel. This requires choosing a connection on the tangent
bundle. A connection can be viewed in various ways. In this chapter,
we shall be using two viewpoints. The first is that a connection 
 provides a means for parallel
transporting vectors along paths. To say it in another way, it
provides an isomorphism of the local structures at different points on
the manifold although the isomorphism may depend on the path. The
second viewpoint is that it provides a method for differentiating
vectors. We shall use $D$ to denote the covariant differentiation
operator.  This discussion works equally for any vector bundle. 
 As part of the definition of a connection, the connection must
 satisfy the Leibiniz rule for tensor products, i.e.\ $D(v \otimes w)
 \ensuremath{:=} Dv \otimes w + v \otimes Dw$ where $v$ and $w$ are sections of
 two vector bundles.\footnote{A 
section of $E$
is a map \ensuremath{\psi}\ from $\ensuremath{\mathcal{Q}}$ into $E$ such that $\ensuremath{\psi}(\ensuremath{q}) \in
E_{\ensuremath{q}}$, i.e. it maps a point \ensuremath{q}\  of \ensuremath{\mathcal{Q}}\ into elements of the
vector fiber over \ensuremath{q}.} It should also commute with contraction, e.g.\ 
if $\omega$ is a section of the dual bundle in which $v$
 is a section of, then $D[\omega(v)] = D[\omega] (v) + \omega(D[v])$. 
In particular, once you have defined a connection on a vector bundle,
then there is an induced connection on its dual space and all of the
various tensor bundles that one can form from these spaces. 


For a metric, 
the connection which is chosen is the Levi-Civita
connection. It is uniquely specified by demanding that the metric is
parallel and that the torsion of the connection vanishes. The fact
that the metric vanishes leads to the commutation of differentiation
and contraction with
either $g$ or $G$. 
For example, if $v$ and $w$ are vector fields, then $D[g(v,w)] =
g(D[v],w) +G(v,D[w])$.

 
For more
general vector bundles, one can also have an inner product on the
fibers. One almost always has a connection on the bundle and demands
that the inner product on the fibers is parallel. This does not define
the connection uniquely. Therefore, we shall call such a bundle a
cc-Hermitian vector bundle where ``cc'' stands for connection-carrying
indicating that a connection has been chosen and the inner product
will always be assumed to parallel under that connection. We shall
generally use $\ensuremath{({\cdot}, {\cdot})}$ to denote the inner product. We
will always use $D$ for the covariant derivative regardless of which
bundles are being discussed. 

We shall now get into the heart of what we need. Given a section \ensuremath{\psi}\ 
of a
cc-Hermitian vector bundle $E$ over the Riemannian manifold \ensuremath{\mathcal{Q}},
we have that $D\ensuremath{\psi}$ is an element of $E \otimes \ensuremath{T^{*} \ensuremath{\mathcal{Q}}}$ which is often
called an $E$-valued 1-form. We might also call this a 1-form valued
section of $E$ which, although literally incorrect, is often useful
terminolgy. The gradient of $\ensuremath{\psi}$ is $\nabla \ensuremath{\psi} \ensuremath{:=} G^{\sharp} (D
\ensuremath{\psi})$. It is an element of $E \otimes \ensuremath{T \ensuremath{\mathcal{Q}}}$ and we might call it a
tangent vector valued section of $E$. For Bohmian mechanics, the
velocity field is a multiple of $\textrm{Im} \ensuremath{({\ensuremath{\psi}}, {\nabla \ensuremath{\psi}})}$ which is
  indeed a section of the tangent bundle. 

The Laplacian of $\ensuremath{\psi}$ is defined to be the metric trace of the second
covariant derivative of $\ensuremath{\psi}$. In symbols, we have $\ensuremath{\Delta} \ensuremath{\psi} \ensuremath{:=} 
G(D^2 \ensuremath{\psi})$.  Note that the second covariant derivative requires the
use of the connection on the cotangent bundle induced by the metric.
This is the divergence of the gradient. Indeed, given a
tangent vector valued section, one can take the contraction of the
derivative with the tanget vector value. In symbolds, $D \cdot \phi
\ensuremath{:=} D_{\ensuremath{j}} \phi^{\ensuremath{j}}$ which one could also view as $\nabla
\cdot \phi \ensuremath{:=} g(\nabla \phi)$. Thus, the divergence of the gradient
is 
\[
\nabla \cdot (\nabla \ensuremath{\psi}) = g(G^{\sharp} D (G^{\sharp}( D \ensuremath{\psi}))) =
g(G^{\sharp} (G^{\sharp} (D^2 \ensuremath{\psi}))) = G(D^2 \ensuremath{\psi}) = \ensuremath{\Delta} \ensuremath{\psi} .
\]
It is a fact that this is equivalent to the Laplace-Beltrami operator
acting on sections of a vector bundle \cite{EL83}. The
Laplace-Beltrami operator is defined on $E$-valued forms by using
Hodge-$*$ duality and the exterior derivative. In the case of a
section, this agrees with the Laplacian described above. For other
forms, it does not. They differ by some curvature terms.  

As a final note, many of the very useful computational tools of
Euclidean analysis do carry over into Riemannian manifolds and
bundles. An important difference is in the coordinate expressions
although one can choose coordinates so that at one particular point,
the expressions look Euclidean. In general, if
$\{e_i\}$ is a local basis for $E$ and $\{dx^l\}$ is a local basis for
\ensuremath{T^{*} \ensuremath{\mathcal{Q}}}, then  the Laplacian of $\ensuremath{\psi} \ensuremath{:=} \ensuremath{\psi}^i e_i$ is
\[
\ensuremath{\Delta} \ensuremath{\psi} = G^{\ensuremath{j} \ensuremath{k}} (\frac{\partial^2 \ensuremath{\psi}^i}{\partial{x^{\ensuremath{j}}}
  \partial{x^{\ensuremath{k}}}} e_i + \ensuremath{\psi}^i D^2 [e_i](\partial_{\ensuremath{j}}
  \otimes \partial_{\ensuremath{k}})  +  \frac{\partial \ensuremath{\psi}^i}{\partial{x^{\ensuremath{j}}}}
   D_{\partial_{\ensuremath{k}}} [e_i] +\frac{\partial \ensuremath{\psi}^i}{\partial{x^{\ensuremath{k}}}}
   D_{\partial_{\ensuremath{j}}} [e_i] 
+ \frac{\partial \ensuremath{\psi}^i}{\partial{x^\ensuremath{j}}}
  e_i 
D_{\partial_{\ensuremath{k}}} [dx^{\ensuremath{j}}] .      
\]


\section{Bohmian dynamics on manifolds} 
A Bohmian dynamics exists in more general situations than the one
described in Chapter \ref{c:bohmint}. In this section, the wave functions will
still be complex-valued functions over the configuration space. 
In fact, we can 
let the configuration space, \ensuremath{\mathcal{Q}}, be an arbitrary Riemannian manifold.
Given a 
metric $g$ on a manifold, the Laplacian may be defined to be
the metric trace of the  
second derivative as explained above. 
The    
 gradient $\nabla \ensuremath{\psi}$  is the (complex) tangent vector
metrically equivalent  to  
 the 1-form $D\ensuremath{\psi}$. Then  $v^{\ensuremath{\psi}}$, the Bohmian velocity field associated to a
wave function\ $\ensuremath{\psi}$,  is  
\begin{equation} \label{ie:rbv}
v^{\ensuremath{\psi}} \ensuremath{:=} \ensuremath{\hbar}\, \textrm{Im} \frac{\ensuremath{({\ensuremath{\psi}}, {\nabla \ensuremath{\psi}})}}{\ensuremath{({\ensuremath{\psi}}, {\ensuremath{\psi}})}}
\end{equation} 
where we have absorbed the mass into the metric.
The time evolution of 
the state $(Q_t, \ensuremath{\psi}_t)$ is given by
\begin{subequations} \label{ie:re}
\begin{align} 
\ensuremath{\frac{d{Q_t}}{d{t}}} =& v^{\ensuremath{\psi}_t}(Q_t)\label{ie:rbe} \\
i \ensuremath{\hbar} \ensuremath{\frac{\partial {\ensuremath{\psi}_t}}{\partial {t}}}  = &- \frac{\ensuremath{\hbar}^2}{2} \ensuremath{\Delta} \ensuremath{\psi}_t  + V \ensuremath{\psi}_t
 \ensuremath{=:} H \ensuremath{\psi}_t.\label{ie:rse}
\end{align}
\end{subequations}
Thus, given $(\ensuremath{\mathcal{Q}}, g, V)$, we have formally specified 
a Bohmian theory.  

The above is suitable for any Riemannian manifold. A primary example
 is
 when the configuration space describes multiple particles moving 
in physical
space. Let the Riemannian manifold $M$, with
metric $g$, be the physical space. Then the configuration space for $N$
distinguished particles moving in $M$ is 
$\ensuremath{\mathcal{Q}} \ensuremath{:=} M^{N}$. 
Let the masses of
the particles be $m_i$. Then the relevant metric on $M^N$, acting
on the tangent space $T_{(q_1, \ldots q_N)} M^N =
 \bigoplus_{i=1}^{N} T_{q_i} M$, is  
\begin{equation*} 
g^N(v_1 \oplus \cdots \oplus v_N, w_1 \oplus \cdots \oplus w_N )
\ensuremath{:=} \sum_{i=1}^N m_i g(v_i,w_i).
\end{equation*} 
Using $g^N$ for the metric,  the Bohmian velocity field 
is given by  \eqref{ie:rbv} and the system evolves according to 
 \eqref{ie:re}.  In Chapter \ref{c:manpart}, we give a more
interesting discussion of many particle systems.
 
\section{The Hamilton-Jacobi equation on manifolds} \label{s:hjman}

As there is no analogue of this for theories whose wave functions are sections
of a bundle, we pause here to establish the standard fact that the modified
Hamilton-Jacobi equation is  a consequence of the Schr\"odinger\ evolution.
Let $\ensuremath{\psi} \ensuremath{=:} R e^{i\frac{S}{\hbar}}$ where $R$ is positive and $S$
is real. Then, the Schr\"odinger\ equation
becomes, after dividing by $e^{i\frac{S}{\hbar}}$, 
\[
i \hbar \ensuremath{\frac{\partial {R}}{\partial {t}}} - R \ensuremath{\frac{\partial {S}}{\partial {t}}} 
= -\frac{\hbar^2}{2}( \ensuremath{\Delta} R + \frac{2 i}{\hbar} \nabla R \cdot \nabla
S + \frac{i}{\hbar} \ensuremath{\Delta} S - 
\frac{1}{\hbar^2} \nabla S \cdot \nabla S) + V R.
\]

Taking the real part of that equation and dividing by $-R$, we have 
\[
\ensuremath{\frac{\partial {S}}{\partial {t}}} = -\frac{1}{2} \nabla S \cdot \nabla S - V 
+\frac{\hbar^2 \ensuremath{\Delta} R}{2R}   
\]
which is the Hamilton-Jacobi equation for the potential $V -
\frac{\hbar^2 \ensuremath{\Delta} R}{2R}$ on the Riemannian manifold \ensuremath{\mathcal{Q}}\ with 
 metric $g$.    

\section{Bohmian dynamics involving bundles} \label{s:bbun}

More generally, we can consider a Bohmian dynamics for wave functions taking
values in some complex vector bundle $E$ over the Riemannian
configuration space \ensuremath{\mathcal{Q}}. 


Our bundles will always be cc-Hermitian vector bundles.
The ``wave functions''
then become smooth sections of the vector bundle; $C^{\infty}
(\ensuremath{\mathcal{Q}}, E)$ will denote the smooth sections while  $C^{\infty}_0
(\ensuremath{\mathcal{Q}}, E)$ will denote the smooth sections with compact
support. 
The local inner
product acting on the vectors $v_q, w_q$ will be denoted by
\ensuremath{({v_q}, {w_q})}. The global inner product on the Hilbert space of wave function
s is the local inner product integrated against the 
Riemannian volume element associated with the metric
$g$; we denote the Hilbert space by $\ensuremath{L^{2}}(\ensuremath{\mathcal{Q}}, E)$.


The potential $V$ is now a 
self-adjoint section of the matrix, or more properly
endomorphism, bundle 
acting on the 
vector bundle's fibers.\footnote{For us, the terms matrix and 
endomorphism will be synonymous; the use of the term matrix does not
imply that we are using a particular basis. We avoid the term linear
operator as that may be confused with an operator acting on the 
Hilbert
space of wave functions. } In the context of a theory involving just spin, 
$V$ is restricted to being a pointwise linear combination of the spin
matrices and 
the identity;  this does exclude a variety of 
 potentials in the case of many particles and generally in the case of
 1 particle where spin-$\frac{1}{2}$ happens to be an exception. 

The equations are the same. Indeed,  $v^{\ensuremath{\psi}}$, 
the Bohmian velocity field associated to a
wave function\ $\ensuremath{\psi}$,  is
\begin{equation} \label{ie:vbv}
v^{\ensuremath{\psi}} \ensuremath{:=} \ensuremath{\hbar}\, \textrm{Im} \frac{\ensuremath{({\ensuremath{\psi}}, {\nabla \ensuremath{\psi}})}}{\ensuremath{({\ensuremath{\psi}}, {\ensuremath{\psi}})}}.
\end{equation}
The time evolution of 
the state $(Q_t, \ensuremath{\psi}_t)$ is given by
\begin{subequations} \label{ie:ve}
\begin{align} 
\ensuremath{\frac{d{Q_t}}{d{t}}} =& v^{\ensuremath{\psi}_t}(Q_t) \label{ie:vbe} \\
i \ensuremath{\hbar} \ensuremath{\frac{\partial {\ensuremath{\psi}_t}}{\partial {t}}} = &- \frac{\ensuremath{\hbar}^2}{2} \ensuremath{\Delta} \ensuremath{\psi}_t  + V \ensuremath{\psi}_t
 \ensuremath{:=} H \ensuremath{\psi}_t. \label{ie:vse}
\end{align}
\end{subequations}


The wave function\ \ensuremath{\psi}\ will always be assumed to be in \ensuremath{L^{2}}\ and  smooth, 
i.e.\ $\ensuremath{\psi}_t \in
C^{\infty}(\ensuremath{\mathcal{Q}}, E) \cap \ensuremath{L^{2}}(\ensuremath{\mathcal{Q}}, E)$.
We shall generally assume that the formal
Hamiltonian, $H$, is essentially 
self-adjoint on the domain of compactly-supported, smooth wave function
s.\footnote{Appendix \ref{a:functan}
 gives an extremely brief discussion about
situations in which 
the Hamiltonian is not essentially self-adjoint on
the domain $C_0^{\infty} (\ensuremath{\mathcal{Q}}, E)$. 
As a particular example, when \ensuremath{\mathcal{Q}}\ is \ensuremath{\mathbb{R}^{3}}\ with
the $z$-axis  removed, the
Laplacian is not essentially self-adjoint on 
$C_0^{\infty}(\ensuremath{\mathcal{Q}}, \ensuremath{\mathbb{C}})$ \cite{AT98}. A choice of a self-adjoint
extension is required corresponding to boundary conditions on the z-axis.


} We shall often use $-\ensuremath{\Delta} +V$ to indicate that we are discussing
the formal Schr\"odinger-type Hamiltonians.
It has been shown that
$-\ensuremath{\Delta} +V$ is essentially self-adjoint on the domain $\ensuremath{C^{\infty}_0}(\ensuremath{\mathcal{Q}}, E)$
 for cc-Hermitian\ vector bundles whose base space
is a complete Riemannian manifold with suitable restrictions on the 
potential $V$;
see \cite{Sh01, Le00} for details. 

\section{The continuity equations for bundle-valued Bohmian dynamics}


We shall now establish the continuity equation for the Bohmian dynamics
descrubed by \eqref{ie:ve}. This is just the standard argument.
The density $\ensuremath{\rho}$ is actually $\ensuremath{\lvert {\ensuremath{\psi}} \rvert}^2 \nu$ where $\nu$ is the
Riemannian volume element. The probability current $\ensuremath{\rho} v^{\ensuremath{\psi}}$ may then be
considered as a $(n-1)$-form and the divergence is the exterior
derivative of the current. In terms of establishing the continuity
equation, 
one can also view $\ensuremath{\rho}$ as a function and then $\ensuremath{\rho} v^{\ensuremath{\psi}}$ would be
a vector field. The divergence is then the contraction of the
derivative of the vector
field. In any event, we want to show that      
\[
\ensuremath{\frac{\partial {\ensuremath{\rho}}}{\partial {t}}} + \nabla \cdot (\ensuremath{\rho} v^{ \ensuremath{\psi}}) = 0.
\]


As a starting point, we can compute from Schr\"odinger's equation that
\[
\begin{split}
\ensuremath{\frac{\partial {\ensuremath{\rho}}}{\partial {t}}} =& \ensuremath{({\ensuremath{\frac{\partial {\ensuremath{\psi}}}{\partial {t}}}}, {\ensuremath{\psi}})} +
\ensuremath{({\ensuremath{\psi}}, {\ensuremath{\frac{\partial {\ensuremath{\psi}}}{\partial {t}}}})} 
\\  = & 
\ensuremath{({\frac{1}{i\ensuremath{\hbar}}(-\frac{\ensuremath{\hbar}^2}{2} \ensuremath{\Delta} \ensuremath{\psi} + V \ensuremath{\psi})}, {\ensuremath{\psi}})} +   
\ensuremath{({\ensuremath{\psi}}, {\frac{1}{i\ensuremath{\hbar}}(-\frac{\ensuremath{\hbar}^2}{2} \ensuremath{\Delta} \ensuremath{\psi} + V \ensuremath{\psi})})}
\\ = &
 \frac{1}{i\ensuremath{\hbar}} \{\frac{\ensuremath{\hbar}^2}{2}(\ensuremath{({\ensuremath{\Delta} \ensuremath{\psi}}, {\ensuremath{\psi}})} - \ensuremath{({\ensuremath{\psi}}, {\ensuremath{\Delta}
 \ensuremath{\psi}})}) + (-\ensuremath{({V \ensuremath{\psi}}, {\ensuremath{\psi}})} + \ensuremath{({\ensuremath{\psi}}, {V \ensuremath{\psi}})})\}.
\end{split}
\]

The potential term automatically vanishes since $V$ is
a self-adjoint matrix-valued function. 
The Laplacian terms may be compactly written as
 $-\ensuremath{\hbar} \textrm{Im} \ensuremath{({\ensuremath{\psi}}, {\ensuremath{\Delta}
\ensuremath{\psi}})} = -\ensuremath{\hbar} \textrm{Im}\{\nabla \cdot \ensuremath{({\ensuremath{\psi}}, {\nabla \ensuremath{\psi}})} - g(\ensuremath{({\nabla
  \ensuremath{\psi}}, {\nabla \ensuremath{\psi}})}) $. Since the last term does not contribute to the
imaginary part, we are  left with the divergence of the probability
current.  


\section{The Pauli equation}


Let \ensuremath{\mathcal{Q}}\  be a Riemannian manifold, with metric $g$, and let $E$
be a cc-Hermitian vector bundle. Physics often requires an additional
structure. This strucutre may be viewed as an entanglement of the
spatial structure with the bundle structure; it comes in the form of a
representation. 

\subsection{A bit of representation theory}
Recall that  a Lie algebra 
representation is a linear map from
a Lie algebra into the set of matrices acting on a vector space
satisfying that the image of a Lie bracket is the commutator of the
images. For a Riemannian manifold, the relevant Lie group is $SO(g)$,
the matrices that preserve the inner product; its 
 Lie algebra is $\ensuremath{\mathfrak{so}(g)}$. We also note that $Spin(g)$ is the
double cover of $SO(g)$ and has the same Lie algebra since the Lie
algebra is the tangent space to the identiy. 

We will assume that our bundle and manifold come equipped with a
Lie algebra representation of $\ensuremath{\mathfrak{so}(g)}$ induced by a representation of
either $SO(g)$ or $Spin(g)$. We shall denote the Lie algebra 
representation as $S: \ensuremath{\mathfrak{so}(g)} \to E\otimes \ensuremath{{E}^{*}}$.  A representation of
a Lie algebra is
a linear map which  preserves Lie brackets. We shall assume that the
Lie group representation is a unitary representation; this does  imply that $S$
takes values in the anti-self-adjoint valued matrices.

We now explain the very standard statement that 
$\ensuremath{\bigwedge^{2}{\ensuremath{T^{*} \ensuremath{\mathcal{Q}}}}}$ is canonically isomorphic, up to a choice of sign, to 
the Lie algebra \ensuremath{\mathfrak{so}(g)}. 
Indeed, the Lie group $SO(g)$ is the set of matrices
$A \in \ensuremath{T^{*} \ensuremath{\mathcal{Q}}} \otimes \ensuremath{T \ensuremath{\mathcal{Q}}}$ which preserve
the metric $g$. That is to say, if $A \in SO(g)$, then for every 
$v, w \in \ensuremath{T \ensuremath{\mathcal{Q}}}$, we have the defining property 
\begin{equation} \label{e3:son}
g(Av, Aw) =
g(v,w).
\end{equation}
This group of matrices is a manifold and the 
Lie algebra is the tangent space to the identity operator. 
More to the point, if we take a path $A(t) \in SO(g)$ 
such that $A(0) = \text{Id}$,  then $\ensuremath{\frac{\partial {A}}{\partial {t}}}$ defines an element of the
Lie algebra;  every such element of the Lie algebra arises in this
way. 
By computing
the derivative of \eqref{e3:son}, we find $
g(\dot A v, w) + g(v, \dot A w) = 0$ which, using the symmetry of $g$,
leads to 
\[
g(\dot A v, w) = -g(\dot A w, v).
\]
The isomorphism between the 2-forms and the Lie algebra is given by
 the map $\dot A \mapsto \omega$ defined by
 $\omega (v, w) \ensuremath{:=}  g(\dot A v, w)$.  By the above
equation, $\omega \in \ensuremath{\bigwedge^{2}{\ensuremath{T^{*} \ensuremath{\mathcal{Q}}}}}$. By using the non-degeneracy of
$g$, we can make a bijective correspondence to the set of antisymmetric
matrices. Thus, given a 2-form $\omega$, one can map it to an antisymmetric
 matrix $A$ and then exponentiate $tA$ to obtain a path $\alpha$ in
 the Lie group such that $\dot{\alpha}(0) = A$. 
 

We may therefore view an element of the Lie algebra \ensuremath{\mathfrak{so}(g)}\ as an element of 
$\ensuremath{\bigwedge^{2}{\ensuremath{T^{*} \ensuremath{\mathcal{Q}}}}}$. The dual
of this space is $\ensuremath{\bigwedge^{2}{\ensuremath{T \ensuremath{\mathcal{Q}}}}}$. We conclude that a
representation may be viewed as a special element of
$\ensuremath{\bigwedge^{2}{\ensuremath{T \ensuremath{\mathcal{Q}}}}} \otimes \ensuremath{{\ensuremath{E}}^{*}} \otimes \ensuremath{E}$.  

\subsection{A generalization of the Laplacian}
 
We start with the already stated observation that 
\[
\ensuremath{\Delta} \ensuremath{\psi} \ensuremath{:=} G \circ (D^2 \ensuremath{\psi}).\] 
It is critical to appreciate this statement. Accepting it almost 
immediately leads to the Pauli equation as an obvious extension. 
Just
to be clear,  $D^2 \ensuremath{\psi} \in (T^*M)^{\otimes
2} \otimes \ensuremath{E}$ is the second covariant derivative, i.e.\ it is a
$(0,2)$ $E$-valued tensor.  The symmetric 
$(2,0)$ tensor $G \in \ensuremath{T \ensuremath{\mathcal{Q}}}^{\otimes 2}$ 
is the dual of the metric $g \in (T^*M)^{\otimes
2}$. We may therefore plug the tensor $G$ into
the tensor $D^2 \ensuremath{\psi}$ to obtain a section of $E$. Note that in the
definition of the Laplacian, we actually should say that $G$ means
 $G \otimes \ensuremath{\text{Id}}$. 

There is an obvious way to modify this definition. We can contract the
second covariant derivative with any matrix-valued $(2,0)$ tensor. 
That is to say,
let $T \in  \ensuremath{T \ensuremath{\mathcal{Q}}}^{\otimes 2} \otimes \ensuremath{E} \otimes \ensuremath{{\ensuremath{E}}^{*}}$. Then $T
\circ D^2 \ensuremath{\psi}$ makes sense and defines the $T$-Laplacian which shall
be denoted by $\ensuremath{\Delta}_T$.


What should $T$ be? 
Since $G$ is
part of the given local information, 
we may assume that the symmetric
part of $T$ 
is $G \otimes \ensuremath{\text{Id}} \in \ensuremath{\bigodot^{2}{\ensuremath{T \ensuremath{\mathcal{Q}}}}} \otimes \ensuremath{E} \otimes \ensuremath{{\ensuremath{E}}^{*}}$. 
This choice is  simple, does
not lead to any new objects, and is certainly invariant under
transformations preserving the metric. This leaves us with the task of
specifying the antisymmetric part of $T$. The easiest possibility, as
originally done, is to take the antisymmetric part to be zero. Assume that we
do not do this. Since the antisymmetric part is an element of 
 $\ensuremath{\bigwedge^{2}{\ensuremath{T \ensuremath{\mathcal{Q}}}}} \otimes \ensuremath{E} \otimes
\ensuremath{{\ensuremath{E}}^{*}}$, we see that the representation $S$ is the same kind of
tensor as the antisymmetric part of $T$. The same reasoning that leads
to choosing $G$ leads us to choosing $S$.


To understand what happens with $T$, we need to understand the second
covariant derivative of a section $\ensuremath{\psi}$. 
The second derivative of $\ensuremath{\psi}$ decomposes into a symmetric part and an
antisymmetric part. The symmetric part is what it is; this is the term
that contributes to the usual Laplacian and is the only part that $G$ notices.
The antisymmetric part is the curvature of the
connection applied to $\ensuremath{\psi}$; this is what $S$ will contract with. 
It is a tensor in $\ensuremath{\psi}$ meaning that it
depends only on $\ensuremath{\psi}$ through its value at that point. It turns out
that if a connection preserves a Hermitian inner product, then its
derivative is anti-self-adjoint in its action on the bundle. See
Appendix \ref{a:difgeom} for an explanation of this fact. We shall use
$F$ to denote the curvature and
 $F$ is an element of $\ensuremath{\bigwedge^{2}{\ensuremath{T^{*} \ensuremath{\mathcal{Q}}}}} \otimes E \otimes
\ensuremath{{E}^{*}}$. 

Let us look at what we have. Taking $\ensuremath{\Delta} \ensuremath{\psi}$ to mean
the metric trace of the second covariant derivative, we see that 
\[
\ensuremath{\Delta}_T \ensuremath{\psi} = \ensuremath{\Delta} \ensuremath{\psi} + \ensuremath{S}(F(\ensuremath{\psi})) 
\]
where $\ensuremath{S}(F)$ is a contraction which leaves a self-adjoint
matrix. Thus, we can write 
\[
\ensuremath{\Delta}_T \ensuremath{\psi} = \ensuremath{\Delta} \ensuremath{\psi} + V \ensuremath{\psi}
\]
for some self-adjoitn $V$ depending on the representation and the
curvature. We do end up with a Pauli-type equation. We also see that
the Bohmian motion is still equivariant with respect to $\ensuremath{\Delta}_T$ for
this choice of $T$. We also mention that one can easily multiply both
$G$ and $S$, independently, in the definition of $T$.  

As an illuminating example, let $\ensuremath{\mathcal{Q}}  = \ensuremath{\mathbb{R}^{3}}$ and $E = \ensuremath{\mathbb{R}^{3}} \times \ensuremath{\mathbb{C}}^2$. We have as a background connection, the trivial
connection denoted by $D$. We are considering the case of 
a spin-$\frac{1}{2}$ particle. Let $S$ be the representation and $\boldsymbol{S}$ its
vector version using the duality between
$(1,0)$ tensors  and antisymmetric $(2,0)$ tensors 
 special to three
dimensions. In particular, $\boldsymbol{S}_x = \frac{\ensuremath{\hbar}}{2} \sigma $ 
is what comes from $S(y\wedge z)$. 
 The curvature is the magnetic field. It is most naturally
a 2-form, but it becomes a vector under duality. The usual Paulie
equation is 
\[
i \ensuremath{\hbar} \ensuremath{\frac{\partial {\ensuremath{\psi}}}{\partial {t}}} 
= -\frac{\ensuremath{\hbar}^2}{2} (\nabla +i\boldsymbol{A})\cdot (\nabla +i\boldsymbol{A}) 
 \ensuremath{\psi} + \mu \boldsymbol{S} \cdot \boldsymbol{B} \ensuremath{\psi}
\]
for some constant $\mu$ and vector potential $\boldsymbol{A}$. The way it
would be written above is the following. Let $\omega$ be the 1-form
equivalent to $A$ under the metric and $\beta$ the 2-form equivalent
to $B$; the 2-form is equal to the exterior derivative of
$\omega$. Then the $T$-based Laplacian equation is 
\[
i \ensuremath{\hbar} \ensuremath{\frac{\partial {\ensuremath{\psi}}}{\partial {t}}} = -\frac{\ensuremath{\hbar}^2}{2} (G + S)(D+i\omega)^2 \ensuremath{\psi}
                   \]
which is the same as the Pauli equation after converting the objects
and expanding. 


\section{Other equivariant motions?}


We shall now ask  whether $T$ leads to a different, yet equivariant, 
 motion for the particle.  The answer is yes. Finding other
 equivariant and invariant motions is not new; see \cite{DG98a, BH95}.
This section is a bit technical and can be skipped without loss of
information. 

Let $T \in \ensuremath{T \ensuremath{\mathcal{Q}}} \otimes \ensuremath{{E}^{*}} \otimes \ensuremath{T \ensuremath{\mathcal{Q}}}  \otimes E$. Then, viewing
$D\ensuremath{\psi} \in \ensuremath{T^{*} \ensuremath{\mathcal{Q}}} \otimes  E$, we can contract on the first two
tensor factors 
of $T$ with $D \ensuremath{\psi}$ to obtain $T \circ D\ensuremath{\psi} \in \ensuremath{T \ensuremath{\mathcal{Q}}} \otimes E$. We
then use the inner product  to obtain a complex tangent vector. Taking
the imaginary part leaves us with a real tangent vector.  Multiplying
by \ensuremath{\hbar}\ and dividing by $\ensuremath{\rho} \ensuremath{:=} {\ensuremath{({\ensuremath{\psi}}, {\ensuremath{\psi}})}}$ leads to a
velocity field, $v^{T, \ensuremath{\psi}}$. 


When do $\ensuremath{\rho}$ and
$v^{T, \ensuremath{\psi}}$ lead to an equivariant motion? In symbols, we are
defining 
 \[
v^{T, \ensuremath{\psi}} \ensuremath{:=} \ensuremath{\hbar} \textrm{Im} \frac{\ensuremath{({\ensuremath{\psi}}, {T\circ D\ensuremath{\psi}})}}{\ensuremath{({\ensuremath{\psi}}, {\ensuremath{\psi}})}}  
\]
and trying to find out what conditions guarantee that the continuity equation
\[
\ensuremath{\frac{\partial {\ensuremath{\rho}}}{\partial {t}}} + \nabla \cdot (\ensuremath{\rho} v^{T, \ensuremath{\psi}}) = 0
\]
is satisfied? 
 
As a starting point, we can compute from Schr\"odinger's equation that
\[
\begin{split}
\ensuremath{\frac{\partial {\ensuremath{\rho}}}{\partial {t}}} =& \ensuremath{({\ensuremath{\frac{\partial {\ensuremath{\psi}}}{\partial {t}}}}, {\ensuremath{\psi}})} +
\ensuremath{({\ensuremath{\psi}}, {\ensuremath{\frac{\partial {\ensuremath{\psi}}}{\partial {t}}}})} 
\\  = & 
\ensuremath{({\frac{1}{i\ensuremath{\hbar}}(-\frac{\ensuremath{\hbar}^2}{2} \ensuremath{\Delta} \ensuremath{\psi} + V \ensuremath{\psi})}, {\ensuremath{\psi}})} +   
\ensuremath{({\ensuremath{\psi}}, {\frac{1}{i\ensuremath{\hbar}}(-\frac{\ensuremath{\hbar}^2}{2} \ensuremath{\Delta} \ensuremath{\psi} + V \ensuremath{\psi})})}
\\ = &
 \frac{1}{i\ensuremath{\hbar}} (-\frac{\ensuremath{\hbar}^2}{2}(-\ensuremath{({\ensuremath{\Delta} \ensuremath{\psi}}, {\ensuremath{\psi}})} + \ensuremath{({\ensuremath{\psi}}, {\ensuremath{\Delta}
 \ensuremath{\psi}})}) + (-\ensuremath{({V \ensuremath{\psi}}, {\ensuremath{\psi}})} + \ensuremath{({V \ensuremath{\psi}}, {\ensuremath{\psi}})}));
\end{split}
\]
we include the spin terms into the potential $V$.
The potential term automatically vanishes since $V$ is
a self-adjoint matrix-valued function. 
The Laplacian terms may be compactly written as
 $-\ensuremath{\hbar} \textrm{Im} \ensuremath{({\ensuremath{\psi}}, {\ensuremath{\Delta}
\ensuremath{\psi}})}$.  
Explicitly, we need
\[
\nabla \cdot (\ensuremath{\rho} v^{T, \ensuremath{\psi}}) = \ensuremath{\hbar} \textrm{Im} \ensuremath{({\ensuremath{\psi}}, {G\circ D^2 \ensuremath{\psi}})}.
\]
Computing the divergence term, we find
\[
\begin{split}
\nabla \cdot (\ensuremath{\rho} v^{T, \ensuremath{\psi}}) &= 
\ensuremath{\hbar} \textrm{Im} (\nabla \cdot \ensuremath{({\ensuremath{\psi}}, {T\circ D\ensuremath{\psi}})}) \\
&= 
\ensuremath{\hbar} \textrm{Im} (\ensuremath{({D_{\ensuremath{k}}\ensuremath{\psi}}, {T^{\ensuremath{j} \ensuremath{k}} D_{\ensuremath{j}}\ensuremath{\psi}})} +
\ensuremath{({\ensuremath{\psi}}, {(D_{\ensuremath{k}}T^{\ensuremath{j} \ensuremath{k}})  D_{\ensuremath{j}}\ensuremath{\psi}})} + 
\ensuremath{({\ensuremath{\psi}}, {T \circ D^2 \ensuremath{\psi}})})
\end{split}
\]
where we use $A_{\ensuremath{j}}B^{\ensuremath{j}}$ to indicate a contraction over the
indices when there may be ambiguity; the 
last term has essentially an unambiguous contraction. We can think
of $\ensuremath{\psi}$, $D\ensuremath{\psi}$ and the symmetric part of $D^2 \ensuremath{\psi}$ 
as  independent of each other. In particular, 
we may choose some to be zero and the others 
non-zero.
We therefore find that we need
\begin{subequations}
\begin{align}
\textrm{Im} \ensuremath{({D_{\ensuremath{k}}\ensuremath{\psi}}, {T^{\ensuremath{j} \ensuremath{k}} D_{\ensuremath{j}}\ensuremath{\psi}})}
 &= 0 \label{e:ce1}\\
\textrm{Im} \ensuremath{({\ensuremath{\psi}}, {(D_{\ensuremath{k}}T^{\ensuremath{j} \ensuremath{k}})  D_{\ensuremath{j}}\ensuremath{\psi}})}
 & = 0 \label{e:ce2}\\
\textrm{Im} \ensuremath{({\ensuremath{\psi}}, {T \circ D^2 \ensuremath{\psi}})}  &= \textrm{Im} \ensuremath{({\ensuremath{\psi}}, {G\circ D^2 \ensuremath{\psi}})}
\label{e:ce3} .
\end{align}
\end{subequations}
Whenever these equations are satisfied, we have an equivariant
motion. 


These equations  rule out many possibilities, but we shall find that
we do have more possibilities than just the usual Bohmian velocity.
The main equation is \eqref{e:ce1}. The essential trick is that we
have complete freedom over the covariant derivatives of \ensuremath{\psi}. Indeed,
as discussed in Chapter \ref{c:taylor}, we can create a \ensuremath{\psi}\ whose
symmetric covariant derivatives are specified at a single point. In
particular, we can choose the first derivatives anyway we like. Pick a
basis of the tangent space at a point \ensuremath{q}; although this discussion only
holds for \ensuremath{q}, since \ensuremath{q}\ is arbitrary, this is sufficient and
we shall suppress the dependency on \ensuremath{q}. 
For a basis element $x$, we 
choose $D [\ensuremath{\psi}] (x) = \phi$ and $D [\ensuremath{\psi}] (y) = 0$ for all the other basis
elements. With this choice, \eqref{e:ce1} becomes 
\[
\textrm{Im} \ensuremath{({\phi}, {T^{xx} \phi})} =0 
\]
which evidently implies that $T^{xx}$ is self-adjoint as this hold for
all $\phi$.
We now choose $D[ \ensuremath{\psi}] (x) = \phi$, $D[ \ensuremath{\psi}](y) =\tau$ and all of the other
basis directional derivatives to vanish. Then \eqref{e:ce1} becomes
\[
\begin{split}
0 =& \textrm{Im}\{ \ensuremath{({\phi}, {T^{xx} \phi})} + \ensuremath{({\tau}, {T^{xy} \phi})}+
\ensuremath{({\phi}, {T^{yx} \tau})} + \ensuremath{({\tau}, {T^{yy} \tau})}\}
\\=&
  \textrm{Im}\{ \ensuremath{({\tau}, {T^{xy} \phi})}+
\ensuremath{({\phi}, {T^{yx} \tau})}\}.
\end{split}
\] 
If we now take $\tau = \phi$, we complete the argument that the
 symmetric part of $T$ is self-adjoint. If we take $\tau = i\phi$,
 then we end up with
\[
0 =  \textrm{Im}\{-i \ensuremath{({\phi}, {T^{xy} \phi})}+
i\ensuremath{({\phi}, {T^{yx} \phi})}= Re \ensuremath{({\phi}, {T^{xy} - T^{yx}  \phi})}
\]
which indicates that the antisymmetric part of $T$ is anti-self-adjoint.
We now move on to \eqref{e:ce3}. The symmetric part of the second
derivative is an object that we can specify arbitrarily at a
point. Thus, repeating an argument  similar to the above, 
we have that the symmetric part of $T$ is equal to $G
\otimes \ensuremath{\text{Id}}$. The antisymmetric part of the second derivative is an
object that we do not have control over. Indeed, it is the curvature
operator. Therefore, if $T$'s antisymmetric part is anti-self-adjoint,
then it couples with the curvature operator in a way to ensure
self-adjointness. That is to say, if $T_A$ is the antisymmetric part of
$T$ and $F$ is the curvature operator, then  
\[
\ensuremath{({\ensuremath{\psi}}, {T\circ D^2 \ensuremath{\psi}})} = \ensuremath{({\ensuremath{\psi}}, {G \circ D^2 \ensuremath{\psi}})} +
\ensuremath{({\ensuremath{\psi}}, {T_A\circ F \ensuremath{\psi}})}.
\]
As $T_A \circ F$ is self-adjoint, its imaginary part vanishes. 

We conclude that the only choice for $T$ is $G \otimes
\ensuremath{\text{Id}} +S$ where $T_A$ is a linear map of the Lie algebra of
\ensuremath{\mathfrak{so}(g)}\ taking values in the Lie algebra of anti-self adjoint
matrices. It is therefore natural to take $T_A \ensuremath{:=} S$ 
to be the representation
map.  Equation \eqref{e:ce2} will be satisfied if $T$ is
parallel which it is if it is composed of $G$ and $S$.


\subsection{The addition of a curl}

For equivariance, we could always add on a divergence-free vector
field to the probability current. 
This will not change the empirical equivalence with standard
quantum mechanics. One can then ask the question as to
whether there is a divergence free velocity field that arises from the
structures at hand. Indeed there is. 

As we have said, $S$ is a anti-self-adjoint matrix-valued 
antisymmetric (2,0) tensor. If we define
\[ \nu \ensuremath{:=} 
\ensuremath{\hbar} \textrm{Im} \ensuremath{({\ensuremath{\psi}}, {S \ensuremath{\psi}})} = -\ensuremath{\hbar} \textrm{Im} \ensuremath{({S\ensuremath{\psi}}, { \ensuremath{\psi}})},
\] 
then we have an antisymmetric (2,0) tensor. 
Such an object can always be used to define a divergence free vector
field. Indeed, take the derivative of $\nu$ and contract; one of the
two vector slots must be chosen which is really a choice of a sign. 
This leaves a vector field.  More explicitly, we note that 
\[
(D [\nu])^{\ensuremath{j}\cdot}_{\ensuremath{j}} 
=    \frac{\ensuremath{\hbar}}{i} (D [\textrm{Im} \ensuremath{({\ensuremath{\psi}}, {S^{\ensuremath{j} \cdot} \ensuremath{\psi}})}])_{\ensuremath{j}}
= \ensuremath{\hbar} \textrm{Im} \ensuremath{({\ensuremath{\psi}}, {S \circ D \ensuremath{\psi}})}
\]
which is essentially what we had in the last section. Thus, adding this to
the probability current will not change the equivariance. 

But we wish to demonstrate that it is equivariant just because $\nu
\in \bigwedge^2 \ensuremath{T \ensuremath{\mathcal{Q}}}$. Let $v$ be the vector field obtained by
contracting the first derivative of $\nu$ with one of its tangent
vector's components. 
To compute the divergence, take the
derivative of the vector field and contract again. We claim that zero
is the answer. 

There are at least two ways to understand this. The first is to view an element
of $\bigwedge^2 (\ensuremath{T \ensuremath{\mathcal{Q}}})$, which $\nu$ is, as an $(n-2)$-form using the
volume form. Then
compute the exterior derivative  to obtain a $(n-1)$-form and this is
equivalent to a tangent vector field. The divergence is computed by
taking the exterior derivative of the $(n-1)$-form and comparing the
resultant $n$-form to the volume form. Since $d^2 =0$, we have that
the divergence vanishes.

The other way is to use the covariant derivative. The only reason to
consider this method is that we are using covariant derivatives and
sometimes this use conflicts with a differential form approach. But in
this case, we end up with the same result.


It is important to
recall
 that covariant differentiation preserves symmetry
types, i.e.\ $D^2 [\nu]$ is an element of $\bigwedge^2(\ensuremath{T \ensuremath{\mathcal{Q}}}) \otimes \ensuremath{T^{*} \ensuremath{\mathcal{Q}}}
\otimes \ensuremath{T^{*} \ensuremath{\mathcal{Q}}}$.   
By computation, we have
\[
(D^2[f \nu])_{{\ensuremath{j}}{\ensuremath{k}}}^{{\ensuremath{j}}{\ensuremath{k}}} = (D^2
[f])_{{\ensuremath{j}}{\ensuremath{k}}} \nu^{{\ensuremath{j}}{\ensuremath{k}}} + 
(D[f])_{\ensuremath{j}} (D[\nu])_{\ensuremath{k}}^{{\ensuremath{j}}{\ensuremath{k}}} + (D[f])_{\ensuremath{k}}
(D[\nu])_{\ensuremath{j}}^{{\ensuremath{j}}{\ensuremath{k}}} + f (D^{\ensuremath{k}} 
[\nu])^{{\ensuremath{j}}{\ensuremath{k}}}_{{\ensuremath{j}}{\ensuremath{k}}} .
\]
The second derivative of a function for a symmetric connection is
symmetric which implies, in conjunction with the antisymmetry of
$\nu$,  that the first term vanishes. The mixed terms
vanish since we clearly have 
\[
(D[f])_{\ensuremath{k}} (D[\nu])_{\ensuremath{j}}^{{\ensuremath{j}}{\ensuremath{k}}} = -
(D[f])_{\ensuremath{k}} (D[\nu])_{\ensuremath{j}}^{{\ensuremath{k}}{\ensuremath{j}}} =  
- (D[f])_{\ensuremath{j}} (D[\nu])_{\ensuremath{k}}^{{\ensuremath{j}}{\ensuremath{k}}}.\]
 Thus,  multiplication by a function commutes with the
divergence operator on antisymmetric 2-tensors.

Let $e$ and $f$ be two vector fields forming part of a local
basis at a point \ensuremath{q}. Then 
\[
\begin{split}
(D^2 [e \otimes f - f \otimes e])_{{\ensuremath{j}}{\ensuremath{k}}}^{{\ensuremath{j}}{\ensuremath{k}}}  =  &
D^2[e]_{{\ensuremath{j}}f}^{\ensuremath{j}} - D^2[f]_{{\ensuremath{j}}e}^{\ensuremath{j}}  
+ 
D[e]_{\ensuremath{j}}^{\ensuremath{j}}  D[f]_{\ensuremath{k}}^{\ensuremath{k}} -
D[f]_{\ensuremath{j}}^{\ensuremath{j}}  D[e]_{\ensuremath{k}}^{\ensuremath{k}} 
\\ &+ 
D[e]_{\ensuremath{k}}^{\ensuremath{j}}  D[f]_{\ensuremath{j}}^{\ensuremath{k}} -
D[f]_{\ensuremath{k}}^{\ensuremath{j}}  D[e]_{\ensuremath{j}}^{\ensuremath{k}} 
+
D^2[f]_{e{\ensuremath{k}}}^{\ensuremath{k}} - D^2[e]_{f{\ensuremath{k}}}^{\ensuremath{k}} .
\end{split}
\]
Clearly, the middle terms, involving the first derivatives, 
vanish since multiplication is commutative. 
The other terms vanish because essentially we are left with  Ricci
curvature terms but in such a way that they cancel. Specifically, 
\[
(D^2[e]_{{\ensuremath{j}}f}^{\ensuremath{j}} - D^2[e]_{f{\ensuremath{k}}}^{\ensuremath{k}}) 
-
( D^2[f]_{{\ensuremath{j}}e}^{\ensuremath{j}}  - D^{\ensuremath{k}}[f]_{e{\ensuremath{k}}}^{\ensuremath{k}})
= 
Ric(e,f) - Ric(f,e).
\]
As Ricci is a symmetric tensor, we have that the right hand side
vanishes. Thus, the divergence vanishes. 

Therefore, one can add this spin-based 
velocity field to the usual Bohmian velocity and still maintain
equivariance. There would seem to be no way to distinguish between
these theories. As described in \cite{BH95}, the non-relativistic limit
of the Bohm motion for the Dirac equation seems to include this
term. We conclude with remarking that if we chose $G\otimes \ensuremath{\text{Id}} +S$
for computing the velocity field, then we are using the same tensor to
compute the Laplacian as well as the velocity field. In other words,
instead of choosing $G$ and $S$, we can view the whole specification
of the theory as choosing the connection (curvature), the object $T$
($G$ and $S$), and the
inner product. 


\section{The Dirac equation}


In this section, we will argue that the structure of the Dirac
equation is based upon what may be called a  square-root of the
tangent bundle. It is this structure which is suggested by the
Bohmian viewpoint. In order to minimize additional structures on
space-time, we actually require our square-root of the tangent bundle
to be a square-root of the metric. We shall see that this puts a severe
constraint on the choices and leads to the Dirac equation being a
square-root of the Klein-Gordon equation. We shall assume we have a
Lorentz metric on space-time and shall denote it by $g$ and its dual
by $G$. As throughout this chapter, almost all statements are about
the local structure and one should have in mind the tangent space at a
point and the fiber at a point rather than the bundle
although we shall omit the reference to the point.

\subsection{Brief background material}


Let \ensuremath{W}\ be a complex vector space. Its dual space, \ensuremath{{\ensuremath{W}}^{*}}, is the
space of complex-linear functionals acting on \ensuremath{W}. The conjugate
dual of $\ensuremath{W}$ is the space of
conjugate-linear functionals acting on \ensuremath{W}; it is denoted  by 
\ensuremath{{\ensuremath{\overline{\ensuremath{W}}}}^{*}} and it is also the
conjugate space of \ensuremath{{\ensuremath{W}}^{*}}. The
dual of \ensuremath{{\ensuremath{W}}^{*}}\ is canonically 
isomorphic to \ensuremath{W}\ under the mapping 
 $\ensuremath{\psi} \mapsto \ensuremath{\psi}(\omega) \ensuremath{:=}
\omega(\ensuremath{\psi})$ where $\ensuremath{\psi} \in \ensuremath{W}$ and $\omega \in \ensuremath{{\ensuremath{W}}^{*}}$. 
The dual of \ensuremath{{\ensuremath{\overline{\ensuremath{W}}}}^{*}}\ is the conjugate space
\ensuremath{\overline{\ensuremath{W}}}\ of \ensuremath{W}. It is the same as the conjugate dual of \ensuremath{{\ensuremath{W}}^{*}}. It
is canonically isomorphic to \ensuremath{W}\ under the  mapping 
$\ensuremath{\psi} \mapsto \ensuremath{\overline{\ensuremath{\psi}}} (\omega) \ensuremath{:=} \ensuremath{\overline{\omega(\ensuremath{\psi})}}$ 
and under the
multiplication that takes $c \odot \ensuremath{\psi} \mapsto \ensuremath{\overline{c}}\ensuremath{\psi}$. This
multiplication allows the appropriate conjugate-linearity to be
expressed in the tensor products. Heuristically, the conjugate space
corresponds to the other choice of what $i$ could mean.


\subsection{Square-roots of tangent spaces and metrics}

From a Bohmian perspective, we want to have a theory of particles
moving on space-time. We restrict our attention to one particle. 
The idea is to try to define a covariant vector field on
space-time in terms of \ensuremath{\psi}. In the non-relativistic theory, one uses the first
derivative of \ensuremath{\psi}\ to do this. But what if the mapping between wave functions
and velocity fields were ultralocal, i.e.\ did not depend on
derivatives, and yet it is still covariant? 

Such maps do exist; we call such a map a  
 square-root of the tangent space. As a starting point, 
the map should be a sesquilinear-quadratic form
taking values in the tangent space; it should be conjugate linear in the
first slot and 
linear in the second. Abstractly, 
$\ensuremath{\nu} \in \ensuremath{{\ensuremath{\overline{\ensuremath{W}}}}^{*}} \otimes \ensuremath{{\ensuremath{W}}^{*}}
\otimes \ensuremath{\mathbb{C}}\ensuremath{T \ensuremath{\mathcal{Q}}}$ which we few as a linear mapping $\ensuremath{\nu}: \ensuremath{\overline{\ensuremath{W}}} \otimes
\ensuremath{W} \to \ensuremath{\mathbb{C}}\ensuremath{T \ensuremath{\mathcal{Q}}}$.  
Given a $\ensuremath{\psi} \in \ensuremath{W}$, we immediately
have that  
$\ensuremath{\nu}(\ensuremath{\overline{\ensuremath{\psi}}}, \ensuremath{\psi}) \in \ensuremath{\mathbb{C}}\ensuremath{T \ensuremath{\mathcal{Q}}}$. 
 We want a real velocity field and therefore demand
that $\ensuremath{\nu}(\ensuremath{\overline{\ensuremath{\psi}}}, \ensuremath{\psi}) $ is a real tangent vector. Then the
Bohmian trajectories are the integral curves of the velocity field. 
 Note that the
parametrization is irrelevant; all we need are the 
 1-dimensional submanifolds whose tangent
spaces correspond to the span of $\ensuremath{\nu}(\ensuremath{\overline{\ensuremath{\psi}}}, \ensuremath{\psi})$.\footnote{As 
we are working in space-time rather than just
space, we know that the position in the manifold is all that should be
relevant.}
From this we can deduce that both \ensuremath{\nu}\ and \ensuremath{\psi}\ should be
thought of as projective objects. We also demand that \ensuremath{\nu}\ be
non-degenerate in the
sense that 
$\ensuremath{\nu}(\ensuremath{\overline{\ensuremath{\psi}}}, \ensuremath{\psi})  =
0 $ if and only if $\ensuremath{\psi} = 0$. In summary, given $\ensuremath{\nu}$ acting on the
 bundle and a
section $\ensuremath{\psi}$, we  have a velocity field for the
law of motion for the particle. 

Before we call $\ensuremath{\nu}$ a square-root of the tangent space, we also
demand that it be covariant. As part of this, we shall need a
representation $R$ of the Lorentz group acting on the bundle. 
To be a square root of the tangent space, \ensuremath{\nu}\ must satisfy, for
every Lorentz
transformation $\Lambda$ acting on the tangent space and every $\ensuremath{\psi}$,
$\phi \in \ensuremath{W}$.
\[
\Lambda (\ensuremath{\nu}(\ensuremath{\psi}, \phi) ) = \ensuremath{\nu} (R(\Lambda) \ensuremath{\psi}, R(\Lambda) \phi). 
\]
Given a covariant $\ensuremath{\nu}$, we then have a covariant law of motion for
the particle.


To form a law that the wave functions must satisfy, we shall need one
more structure, which we call $\mu$,  an appropriately linear
mapping  $\ensuremath{{\ensuremath{\overline{\ensuremath{W}}}}^{*}} \to \ensuremath{W}$. Thinking of \ensuremath{\nu}\ as a
mapping from $\ensuremath{W} \to \ensuremath{{\ensuremath{\overline{\ensuremath{W}}}}^{*}} \otimes \ensuremath{T \ensuremath{\mathcal{Q}}}$, we see that $\mu \circ
\ensuremath{\nu}$ leads to a mapping $\gamma : \ensuremath{W} \to \ensuremath{W} \otimes
\ensuremath{\mathbb{C}} \ensuremath{T \ensuremath{\mathcal{Q}}}$. Assuming that $\mu$ is invertible, which we do, we can define a,
not necessarily positive definite, 
sesquilinear product $\ensuremath{\langle \ensuremath{\psi}, \ensuremath{\psi} \rangle} \ensuremath{:=} \{\mu^{-1} (\ensuremath{\psi})\}
(\ensuremath{\overline{\ensuremath{\psi}}})$. 
Thus, we have that $\ensuremath{\nu}(\ensuremath{\overline{\ensuremath{\psi}}}, \ensuremath{\psi}) =
\ensuremath{\langle \ensuremath{\overline{\ensuremath{\psi}}}, \gamma \ensuremath{\psi} \rangle}$. As this should be real for all \ensuremath{\psi}\, we
have that $\gamma$ should be appropriately self-adjoint with
respect to $\ensuremath{\langle \cdot, \cdot \rangle}$. We also demand that $\mu$ be
covariant or, equivalently, that 
$\ensuremath{\langle R(\Lambda) \ensuremath{\psi}, R(\Lambda) \phi \rangle} =
\ensuremath{\langle \ensuremath{\psi}, \phi \rangle}$. 

What else can we generate with these structures? With \ensuremath{\langle \cdot, \cdot \rangle}\ and
$\gamma$, we can construct more tensor objects. For example,
$T \ensuremath{:=} \ensuremath{\langle \ensuremath{\overline{\ensuremath{\psi}}}, \gamma^2 \ensuremath{\psi} \rangle}$ 
is a $(2,0)$ tensor. If we wish to minimize
new structures on space-time, then it would seem reasonable to demand
that the symmetric part of $T$ is actually $G
\ensuremath{\langle \ensuremath{\overline{\ensuremath{\psi}}}, \ensuremath{\psi} \rangle}$ which immediately means that
 the symmetric part of $\gamma^2$ is 
$G \otimes \ensuremath{\text{Id}}$. Making
this one  choice leads to very few possibilities,
after taking into account equivalences, of
what the pair $(\ensuremath{\nu}, \mu)$ represent.\footnote{The 
author does not know whether covariance of $\ensuremath{\nu}$ and $\mu$ are
sufficient to deduce this.} Indeed, $\gamma$ is a
representation map for a Clifford algebra as we shall now explain. 


The refrence for this material is \cite{GW98}. 
A Clifford algebra for the vector space $V$, endowed with the metric
$\beta$,  is a an associative algebra with unit 1 equipped with a linear
mapping $\alpha : V \to \ensuremath{\mathrm{Cliff}(V, \beta)}$ such that $\alpha(V)$
generates the algebra and the anticommutator of the images is a multiple of the
identity, i.e.\ $\{\alpha(x), \alpha(y)\} =  \alpha(x) \alpha(y) +
 \alpha(y) \alpha(x) =\beta(x,y) 1$
for all $x, y \in V$.\footnote{One also requires a universality
  property which we shall not state.} It is a fact that a Clifford
algebra is unique up to isomorphism. It is constructed from the tensor
algebra of $V$ by essentially viewing $x \otimes y + y\otimes x$ as
$\beta(x,y) 1$, i.e.\ mod out by the ideal generated by that
equation. If we allow $\beta$ to be identically zero, then the
Clifford algebra generated is the exterior algebra. For us, the vector
space $W$ is $\ensuremath{\mathbb{C}} \ensuremath{T^{*} \ensuremath{\mathcal{Q}}}$ and the inner product $\beta$ is $G$. 

A space of spinors for $(V, \beta)$ is a vector space $W$ equipped with
a linear map $\gamma : V \to W\otimes \ensuremath{{W}^{*}}$ such that the only
invariant subspaces of $W$ under $\gamma(V)$ is $W$ and 0 and that 
$\{ \gamma(x), \gamma(y)\} = \beta(x,y) \ensuremath{\text{Id}}$. The map $\gamma$ then
extends to an irreducible representation of the Clifford algebra and
all irreducible representations arise in this way.
 For us, $\gamma$
is the same $\gamma$ we were discussing before once we convert the
spinor $\gamma$ to a map of the form $W \to V^* \otimes W$. It is a
fact that for even dimensions, there is exactly one space of spinors
while for odd dimensions, there is two inequivalent spaces of spinors.  

 We therefore have a
complete specification of the additional structures in four dimensions
if the Clifford algebra representation is irreducible. Furthermore,
$\gamma^2 = G \otimes \ensuremath{\text{Id}} + \ensuremath{S}$, where $\ensuremath{S}$ is a
particular 
 representation, up to equivalence, of the
Lorentz Lie algebra, \ensuremath{\mathfrak{{so}}}(g). For even dimensions, the Lie algebra
representation decomposes into     the two, inequivalent, half-spin
representations of $\ensuremath{\mathfrak{so}(g)}$. For odd dimensions, the Lie algebra
representation is irreducible. Notice that the covariance plays no
role in determining the uniqueness of these structures; covariance even seems
to be implied from demanding that the symmetric part of $\gamma^2$ 
is the metric.   

The pair $(\ensuremath{\nu}, \mu)$, or equivalently $(\gamma, \ensuremath{\langle \cdot, \cdot \rangle})$,  is
what we call a square-root of the metric since, when we square
$\gamma$, the symmetric part is the metric. 


The usual Gamma matrices also give rise to an irreducible
representation of the Clifford algebra for $G$. By the uniqueness
cited above, we find that a square-root of the Lorentz metric
on four-dimensional space corresponds to an equivalent representation
as the usual Gamma matrices. 

\subsection{Formulating the Dirac equation}

We shall continue to use $D$ for the covariant derivative and it shall
be used for differentiating any object such as vector fields or
sections of the bundle $E$. 


We start by noting that $\ensuremath{\nu}$ and $\mu$ should be
parallel. Equivalently, we demand that, for any section $\ensuremath{\psi}$, we have
\begin{align}
D \ensuremath{\nu}(\ensuremath{\overline{\ensuremath{\psi}}}, \ensuremath{\psi}) &= \ensuremath{\nu} (D \ensuremath{\overline{\ensuremath{\psi}}}, \ensuremath{\psi}) + \ensuremath{\nu}(\ensuremath{\overline{\ensuremath{\psi}}}, D \ensuremath{\psi}) \\
D \ensuremath{\langle \ensuremath{\overline{\ensuremath{\psi}}}, \ensuremath{\psi} \rangle} &= \ensuremath{\langle D \ensuremath{\overline{\ensuremath{\psi}}}, \ensuremath{\psi} \rangle} + \ensuremath{\langle \ensuremath{\overline{\ensuremath{\psi}}}, D\ensuremath{\psi} \rangle}.
\end{align}


Using the irreducibility of the Clifford algebra representation, it
can be shown that the set of possible connections is characterized by
the class of 1-forms rather than general matrix-valued 1-forms.  


To obtain the Dirac equation, one could try to formulate the simplest
equation with the structures at hand. What do we have? We have a
mapping $\gamma$ taking wave functions to elements of $\ensuremath{\mathbb{C}} \ensuremath{T \ensuremath{\mathcal{Q}}} \otimes \ensuremath{W}$. The
covariant derivative takes wave functions to elements of $\ensuremath{\mathbb{C}} \ensuremath{T^{*} \ensuremath{\mathcal{Q}}} \otimes
\ensuremath{W}$. Both are covariant operations. If we contract
with these two operators, then we have an invariant operator taking
sections of $E$ to sections of $E$. 
This is 
the Dirac operator which is defined to be $\ensuremath{ \slash \! \! \!} \partial \ensuremath{\psi} \ensuremath{:=}
\gamma^{\ensuremath{j}} D_{\ensuremath{j}} \ensuremath{\psi}$. Seeking an invariant equation for
the wave function\ using the Dirac operator, we see that the Dirac equation is
an obvious candidate:
\begin{equation}
\frac{\hbar}{i} \ensuremath{ \slash \! \! \!} \partial \ensuremath{\psi} = a \ensuremath{\psi}.
\end{equation}


There is another way to argue for this equation.   
As before, we want to minimize the structures placed on space-time. 
Given a vector field, one can take its derivative. Demanding that this
vanishes leads to a rather uninteresting vector field. If we demand,
on the other hand, that only the divergence vanishes, then a natural
candidate is 
 the Dirac equation. This is the space-time version of
equivariance. 
The velocity field being divergence-free is the equation 
\[
0 = D_{\ensuremath{j}}\ensuremath{\langle \ensuremath{\overline{\ensuremath{\psi}}},  \gamma^{\ensuremath{j}} \ensuremath{\psi} \rangle} = \ensuremath{\langle D_{\ensuremath{j}}
  \ensuremath{\overline{\ensuremath{\psi}}},  
\gamma^{\ensuremath{j}} \ensuremath{\psi} \rangle} + \ensuremath{\langle \ensuremath{\overline{\ensuremath{\psi}}},  \ensuremath{ \slash \! \! \!} \partial
  \ensuremath{\psi} \rangle}
= \ensuremath{\overline{\ensuremath{\langle \ensuremath{\overline{\ensuremath{\psi}}}, \ensuremath{ \slash \! \! \!} \partial \ensuremath{\psi}  \rangle}}} 
+ \ensuremath{\langle \ensuremath{\overline{\ensuremath{\psi}}},  \ensuremath{ \slash \! \! \!} \partial
  \ensuremath{\psi} \rangle}
\]
where we have used the previously mentioned 
fact that $\ensuremath{\langle \ensuremath{\overline{\tau}}, {\gamma \phi} \rangle} =
\ensuremath{\langle \ensuremath{\overline{\gamma \tau}}, \phi \rangle}$ for all $\tau, \phi \in \ensuremath{W}$.
Thus, 
we need the
real part of $\ensuremath{\langle \ensuremath{\overline{\ensuremath{\psi}}},  \ensuremath{ \slash \! \! \!} \partial  \ensuremath{\psi} \rangle}$ to
vanish. This 
certainly follows if  $i b \ensuremath{\psi} = \ensuremath{ \slash \! \! \!} \partial
\ensuremath{\psi}$ for some 
real constant $b$ which is the Dirac equation. 

A final fact is  that the Dirac operator is the square root of a
Pauli-type Laplacian. Indeed, 
\[
a^2 \ensuremath{\psi} = (\frac{\hbar}{i} \ensuremath{ \slash \! \! \!} \partial) 
(\frac{\hbar}{i} \ensuremath{ \slash \! \! \!} \partial) \ensuremath{\psi} 
= - \hbar^2 (\gamma^2)^{{\ensuremath{j}}{\ensuremath{k}}} D^2_{{\ensuremath{j}}{\ensuremath{k}}} \ensuremath{\psi} 
= - \hbar^2 (\ensuremath{\Delta} \ensuremath{\psi} + \ensuremath{S}(
F) \ensuremath{\psi})
\]
where $
F$ is the curvature of the connection. 
Thus, a solution to the Dirac equation also satisfies the
Klein-Gordon-Pauli equation.

\clearpage
\bibliographystyle{abbrv}
\bibliography{references}
\end{document}
